\documentclass[11pt,a4paper]{article}
\usepackage[margin=1in]{geometry}
\usepackage[T1]{fontenc}
\usepackage[utf8]{inputenc}
\usepackage{graphicx}
\usepackage{hyperref}
\usepackage{enumitem}
\usepackage{float}
\usepackage{amsmath}

\title{Advanced Image Processing Methods\\Homework 3 Report}
\author{Peter Dobosi}
\date{\today}

\begin{document}
\maketitle

\section*{Goal}
The task is to investigate how noise reduction and compression parameters affect reconstructed image quality and compression ratio.
The pipeline applies Gaussian denoising at multiple kernel sizes, JPEG compression at varying quality levels, and evaluates results using full-reference image quality metrics (MSE, PSNR, SSIM).

\section*{Method Summary}
\begin{enumerate}[leftmargin=*]
    \item \textbf{Dataset Loading (Task~1)}:
        10 clean reference images (Kodak dataset, \texttt{kodim01}--\texttt{kodim10}) and their noisy counterparts are loaded dynamically using \texttt{Path.glob}.
        The clean images are indexed from 1 while the noisy images are indexed from 0, so pairing is done by sorted order rather than by filename index.
        Each pair is validated for successful loading and matching dimensions.

    \item \textbf{Noise Reduction (Task~2)}:
        The noisy images are corrupted with normally-distributed (Gaussian) noise, as indicated by the \texttt{norm\_dist} suffix in the dataset folder name.
        A Gaussian low-pass filter (\texttt{cv2.GaussianBlur}) is the appropriate choice for this noise type: it optimally attenuates additive Gaussian noise while preserving image structure.
        Three kernel sizes are evaluated: $3\times3$, $5\times5$, and $7\times7$.
        Sigma is derived automatically by OpenCV from the kernel size.
        Every noisy image is denoised with all three kernel sizes.

    \item \textbf{Image Compression (Task~3)}:
        Each denoised image is compressed using WebP lossy encoding (\texttt{cv2.imencode} with \texttt{IMWRITE\_WEBP\_QUALITY}).
        The compressed buffer is immediately decoded back (\texttt{cv2.imdecode}) to obtain the reconstructed image, simulating a lossy round-trip.
        Quality levels from 10 to 100 in steps of 10 are evaluated.
        Every denoised image (3 kernel sizes $\times$ 10 images) is compressed at all 10 quality levels, yielding 300 reconstructed images.
        The compressed buffer size in bytes is recorded for computing compression ratios.

    \item \textbf{Quality Measurement (Task~4)}:
        Each reconstructed image (after denoising + WebP compression) is compared against the corresponding clean reference using three full-reference image quality assessment (FR~IQA) metrics:
        \begin{itemize}[leftmargin=1.5em]
            \item \textbf{MSE} -- Mean Squared Error, computed as the average of the squared per-pixel differences in float64.
            \item \textbf{PSNR} -- Peak Signal-to-Noise Ratio in dB, via \texttt{cv2.PSNR}.
            \item \textbf{SSIM} -- Structural Similarity Index (Wang et al., 2004), computed per-channel using an $11\times11$ Gaussian window ($\sigma=1.5$) and averaged across all pixels and channels.
        \end{itemize}
        The compression ratio is also recorded as $\text{raw\_size} / \text{compressed\_size}$.
        All metrics are computed for every (kernel size, quality) combination, yielding $3 \times 10 = 30$ averaged measurement points.

    \item \textbf{Evaluation (Task~5)}:
        % TODO: Describe plots and optimal parameter selection.
\end{enumerate}

\section*{Filter Selection Rationale}
The Gaussian filter was selected based on two factors:
\begin{itemize}[leftmargin=*]
    \item \textbf{Noise type match}: the dataset folder name (\texttt{dobosi\_peter\_laszlo\_norm\_dist}) indicates normally-distributed noise.
        Gaussian blur is the optimal linear filter for additive Gaussian noise, as it minimizes mean squared error in the frequency domain.
    \item \textbf{Simplicity}: among the candidate filters listed in the assignment (averaging, Gaussian, median, bilateral), the Gaussian filter requires only the kernel size parameter, with sigma derived automatically.
        This makes systematic comparison across kernel sizes straightforward.
\end{itemize}

A larger kernel provides stronger smoothing and better noise suppression, but at the cost of blurring fine image details.
The quality metrics in Task~4 will quantify this trade-off.

\section*{Compression Quality Effect}
Higher WebP quality values produce larger buffers but preserve more image detail.
The table below shows averaged compression ratios (across all 10 images) for each kernel size and selected quality levels:

\begin{center}
\begin{tabular}{l|c c c c c}
    \textbf{Kernel} & \textbf{q=10} & \textbf{q=30} & \textbf{q=50} & \textbf{q=70} & \textbf{q=100} \\
    \hline
    $3\times3$ & 134:1 & 83:1 & 59:1 & 43:1 & 7:1 \\
    $5\times5$ & 153:1 & 99:1 & 73:1 & 55:1 & 9:1 \\
    $7\times7$ & 174:1 & 118:1 & 90:1 & 70:1 & 13:1 \\
\end{tabular}
\end{center}

Larger kernels smooth out more detail, which makes the image more compressible (higher ratios).
At low quality settings the compressed size is very small but lossy artifacts accumulate on top of the denoising loss.
At quality~100, WebP uses near-lossless encoding, resulting in much larger files with minimal additional compression distortion.

\section*{Quality Metrics Summary}
The following table shows the averaged MSE, PSNR, and SSIM for each kernel size and selected quality levels:

\begin{center}
\begin{tabular}{l|c c c|c c c|c c c}
    & \multicolumn{3}{c|}{\textbf{MSE}} & \multicolumn{3}{c|}{\textbf{PSNR (dB)}} & \multicolumn{3}{c}{\textbf{SSIM}} \\
    \textbf{Quality} & $3{\times}3$ & $5{\times}5$ & $7{\times}7$ &
                       $3{\times}3$ & $5{\times}5$ & $7{\times}7$ &
                       $3{\times}3$ & $5{\times}5$ & $7{\times}7$ \\
    \hline
    10  & 159 & 201 & 246 & 26.9 & 26.0 & 25.2 & 0.739 & 0.700 & 0.662 \\
    30  & 123 & 169 & 221 & 28.0 & 26.8 & 25.7 & 0.788 & 0.738 & 0.689 \\
    50  & 106 & 154 & 209 & 28.7 & 27.2 & 26.0 & 0.815 & 0.760 & 0.704 \\
    70  &  97 & 146 & 201 & 29.1 & 27.5 & 26.2 & 0.832 & 0.776 & 0.716 \\
    100 &  81 & 129 & 185 & 30.0 & 28.2 & 26.6 & 0.857 & 0.807 & 0.745 \\
\end{tabular}
\end{center}

The $3\times3$ kernel consistently achieves the best quality across all metrics and quality levels, because it removes the least amount of image detail during denoising.
However, this comes at the cost of less noise suppression and lower compression ratios compared to larger kernels.

\section*{Results}
% TODO: Add figures and discussion after Tasks 3--5 are implemented.

\section*{Reproducibility}
From the \texttt{homework\_3/} directory:
\begin{itemize}[leftmargin=*]
    \item Install dependencies: \texttt{pip install -e .}
    \item Run pipeline: \texttt{python src/mw79on\_submission\_hw3/main.py}
\end{itemize}
The script generates output figures in \texttt{output/}.
This report is provided as \texttt{REPORT.pdf}.

\end{document}
